\subsection{Where the analysis time goes}
\label{sec:anatomy}

\paragraph{Phase 1 costs about 70\,\textmu s per trace line.}
Figure~\ref{fig:phase1} times upstream's Phase~1 on its own (the first 61 lines of the
unmodified \texttt{processing.py}, cut programmatically) against the rewrite's
in-memory version, which writes the same files. On the WSL root file system the
upstream loop settles at 74\,\textmu s per line for long traces: 7.4\,s for
$10^5$ lines and 73.8\,s for $10^6$, against 0.32\,s and 2.4\,s for the rewrite. With
the working directory on ext4, XFS, Btrfs or F2FS it costs 68--79\,\textmu s per line
at $10^5$ lines, for one traced file or 1{,}000, so the file system hardly matters. A
system-call count explains it. For each trace line, Phase~1 makes about 41 calls: 30
one-byte \texttt{read}s, because upstream opens the trace with \texttt{buffering=0}
and then calls \texttt{readlines()}, which on an unbuffered file reads one byte at a
time, and then \texttt{openat}, \texttt{fstat}, an \texttt{ioctl} (the terminal check
of Python's \texttt{open}), five \texttt{lseek}s, one \texttt{write} and one
\texttt{close} to append a single line to \texttt{<inode>.txt}. The fork copied this
code into both of its paths.

\begin{figure}[t]
\centering
\figfile{width=\columnwidth}{../figs/e3_phase1_F1.pdf}
\caption{Phase~1 wall time against trace length, upstream versus rewrite, one traced
file. Medians of ten runs.}
\label{fig:phase1}
\end{figure}

\paragraph{Phase 2 and the rest of the pipeline cost fourteen processes per file.}
Table~\ref{tab:anatomy} (\S\ref{sec:background}) showed the per-file work. On the
real traces of E0, Table~\ref{tab:real} and Figure~\ref{fig:breakdown} split the whole
upstream pipeline into its steps. When a workload touches many files, the per-file steps
dominate. For \texttt{tar}, which read 5{,}100 files on the test file system, the
pipeline takes 165\,s: 97\,s in Phase~2, 56\,s in \texttt{hotness.sh} and 10\,s in
\texttt{merge.py}, against 1.7\,s in Phase~1. The rewrite does all three steps in
1.28\,s (129$\times$, 95\% CI 121--133). For \texttt{fio}'s 1{,}000 files the pipeline
takes 40.4\,s against 0.62\,s (65$\times$). When a workload touches one file, Phase~1
dominates instead: 83\% of the 0.75\,s for \texttt{read\_seq} and 92\% of the 8.6\,s
for \texttt{sqlite3}, whose trace has 160{,}935 lines (69\% of them for its database
file, the only traced file on the test file system). The
best \texttt{--parallel} setting makes
\texttt{processing.py} 7.4$\times$ faster than upstream on \texttt{tar} (13.4\,s against
98.7\,s) and 3.4$\times$ on \texttt{fio}, but slower than upstream on every one-file
trace (1.81\,s against 0.67\,s for \texttt{read\_seq}), because of its one-second floor
(\S\ref{sec:ladder}).

\begin{table*}[t]
\caption{Real traces, warm cache: wall time (s, median of seven runs) of upstream
\texttt{processing.py} (a), the best \texttt{--parallel} setting (d), the rewrite (c),
and of the whole pipeline (\texttt{processing.py}, \texttt{merge.py},
\texttt{hotness.sh}) against the rewrite doing all three steps.}
\label{tab:real}
\small
\centering
\inputtable{../results/tab_real.tex}
\end{table*}

\begin{figure*}[t]
\centering
\figfile{width=0.85\textwidth}{../figs/e2real_breakdown_warm.pdf}
\caption{Left: share of the upstream pipeline's time by step on each real trace. Right:
total time, upstream pipeline against the rewrite (median of seven runs, warm cache).}
\label{fig:breakdown}
\end{figure*}
